%==============================================================================
% Closing the Variance Gap: Bochner PINN for Lévy-Driven SPDEs
% Hubert Lipski
%
% Version 4: incorporates all experimental findings from the Heston-V
% diagnosis campaign (subgrid overfit, randomized subgrid fix, cumulant
% matching for kurtosis, Lévy intensity tau-sweep).
%
% Key changes from v3:
%  - Abstract: Heston-V "structural failure" replaced by diagnosis + fix
%  - §1.3 Contributions: three new bullets (subgrid overfit, cumulant, Lévy tau)
%  - §2 Method: new §2.6 randomized spectral sub-grid
%  - §3.4 Theory Remark: reframed from "structural failure" to "initial diagnosis"
%  - §5 Results: new §5.5 Heston-V diagnosis with 2×4 ablation table
%  - §6 Discussion: "Heston variance channel: diagnosed and fixed" replaces
%   "a documented failure"
%  - §7 Conclusion: updated
%  - Appendix A: Lévy intensity threshold sweep (negative result)
%==============================================================================

\documentclass[review,12pt]{elsarticle}

\usepackage[utf8]{inputenc}
\usepackage[T1]{fontenc}
\usepackage{amsmath, amsthm, amssymb}
\usepackage{booktabs}
\usepackage{multirow}
\usepackage{siunitx}
\usepackage{graphicx}
\usepackage{hyperref}
\usepackage[disable]{microtype}
\usepackage{tikz}
\usetikzlibrary{positioning, arrows.meta, shapes.geometric}

\graphicspath{{figures/}}

\newtheorem{remark}{Remark}
\newtheorem{proposition}{Proposition}
\newtheorem{lemma}{Lemma}
\newtheorem{corollary}{Corollary}
\newtheorem{conjecture}{Conjecture}

\journal{Journal of Computational Physics}

\begin{document}

\title{Supplementary Material:\\ Closing the Variance Gap: Bochner PINN for Physics-Informed Learning on Stationary Random Fields}
\author{Hubert Lipski}
\date{\today}
\maketitle

\noindent This document collects the proof of
Proposition~3.1 (rigorous Gaussian single-realisation convergence
rate under the kernel-smoothed periodogram loss), the supporting
lemmas, the L\'evy partial-result proofs together with a frank
discussion of the obstacles to a fully rigorous L\'evy extension,
the corollaries, the non-separable extension, the NTK conjectural
training-time bound, the L\'evy intensity threshold sweep
appendix, and (S4) a full-law convergence result for L\'evy drivers
in the Zolotarev metric $\zeta_s$ that removes the $4\sqrt J$
Gaussian-projection floor and covers $\alpha$-stable forcing,
$s < \alpha$. Numbering restarts at A.

\appendix

% ─────────────────────────────────────────────────────────────────────────────
\section*{S1 \quad Proof of Proposition~3.1 and its extensions}
\label{app:convergence-proofs}
% ─────────────────────────────────────────────────────────────────────────────

This appendix collects the four-step proof of
Proposition~\ref{thm:main} (main theorem), the supporting
lemmas (log-to-linear control, expected spectral quadrature error,
single-realisation spectral consistency under mixing,
Plancherel--Gelbrich bound, the Mishra--Molinaro adaptation), the
two corollaries (variance-gap rate, L\'evy/non-separable extensions),
the discussion of the rate, and the end-to-end training-time bound
originally drafted as Section~3.3 of the main text. The material is
moved here verbatim; no logical change.

The proof proceeds in four steps.

\begin{lemma}[Log-to-linear control on bounded spectra]
\label{lem:log-linear}
Let \(f, g : \mathcal G_s \to \mathbb R_{>0}\) with
\(\mu_{\min} \le f, g \le \mu_{\max}\). Then
\[
\|f - g\|^2_{L^2(\mathcal G_s)}
\;\le\; (\ln 10)^2\,\mu_{\max}^2\;
\|\log_{10} f - \log_{10} g\|^2_{L^2(\mathcal G_s)}.
\]
\end{lemma}
\begin{proof}
The function \(\varphi(x) = \log_{10} x\) has derivative
\(\varphi'(x) = (x \ln 10)^{-1}
\in [(\mu_{\max}\ln 10)^{-1},\, (\mu_{\min}\ln 10)^{-1}]\)
on \([\mu_{\min}, \mu_{\max}]\). By the mean-value theorem,
\(|f - g| \le \mu_{\max}\ln 10\;|\log_{10} f - \log_{10} g|\)
at each grid point. Squaring and summing gives the claim.
\end{proof}

\begin{remark}[Why we apply Lemma~\ref{lem:log-linear} to the
kernel-smoothed periodogram, not to the raw periodogram]
\label{rem:why-smoothed}
The raw periodogram $\hat S^{\rm xy}_{\mathcal G_s}(u)(k)$ at a
single Fourier mode is distributed as $\tfrac12 S^{\rm xy}(k)\,\chi^2_2$
(an exponential with mean $S^{\rm xy}(k)$), so its support is
$[0,\infty)$ and the upper bound $\mu_{\max}$ of (A3) is exceeded
with positive probability at every mode. Lemma~\ref{lem:log-linear}
therefore cannot be applied to the raw periodogram with the
operator-prescribed $\mu_{\max}$; doing so requires an enlarged
realisation-dependent constant $\mu_{\max}^{(\hat S)}$ that does not
match the assumption (A3). The rigorous result avoids this issue
by applying Lemma~\ref{lem:log-linear} only to the
\emph{kernel-smoothed} periodogram $\hat S^{\rm xy}_h$
(Lemma~\ref{lem:single-realisation}); under (A2) and a kernel
$\kappa$ satisfying the integrability conditions stated there, the
smoothed periodogram is uniformly bounded by a multiple of
$\mu_{\max}$ in expectation and, by the variance bound of
Lemma~\ref{lem:single-realisation}, with high probability up to a
factor controlled by $(N_s h_N^d)^{-1/2}$ which vanishes at the
chosen bandwidth $h_N^\star$.
\end{remark}

Applied to \(f = \widehat S^{\mathrm{xy}}_{h}(u_\theta)\),
\(g = \widehat S^{\mathrm{xy}}_{h}(u_{\mathrm{FD}})\) at the
kernel-smoothed bandwidth $h = h_N^\star$, and combined with
\eqref{eq:tolerance} interpreted on the smoothed loss:
\begin{equation}
\bigl\|\widehat S^{\mathrm{xy}}_h(u_\theta)
- \widehat S^{\mathrm{xy}}_h(u_{\mathrm{FD}})
\bigr\|^2_{L^2}
+ \bigl\|\widehat S^t_h(u_\theta)
- \widehat S^t_h(u_{\mathrm{FD}})
\bigr\|^2_{L^2}
\;\le\; 2(\ln 10)^2\,\mu_{\max}^2\;\varepsilon.
\label{eq:log-linear-applied}
\end{equation}

\begin{lemma}[Expected spectral quadrature error]
\label{lem:quadrature}
Under (A2), for every \(u \in L^2_{\mathbb P}\) stationary with
spatial spectral density
\(S^{\mathrm{xy}} \in H^{s_x}(\widehat\Omega)\),
\[
\bigl\|\mathbb E\bigl[\widehat S^{\mathrm{xy}}_{\mathcal G_s}(u)\bigr]
- S^{\mathrm{xy}}(u)\bigr\|^2_{L^2(\widehat\Omega)}
\;\le\; C_{\mathrm{quad},x}\,(N_s^x)^{-2s_x/d},
\]
with \(C_{\mathrm{quad},x} = C_x\,\|\mu^{\mathrm{xy}}\|^2_{H^{s_x}}\)
depending only on \(\mathrm{diam}(\Omega)\) and \(d\). An analogous
bound holds for the temporal periodogram with rate \((N_s^T)^{-2s_t}\).
\end{lemma}
\begin{proof}
This is the standard \(L^2\) trapezoidal-rule error rate for
Sobolev functions on a \(d\)-dimensional torus
(cf.~\cite[Sec.~6.5]{Trefethen2000}) applied to the expected
periodogram \(\mathbb E[\widehat S^{\mathrm{xy}}_{\mathcal G_s}(u)]\),
which equals a Dirichlet-kernel smoothing of \(S^{\mathrm{xy}}\) and
inherits its \(H^{s_x}\) regularity. The constant absorbs the
\(H^{s_x}\) norm.
\end{proof}

\begin{remark}[Single-realisation variance and the proof gap]
\label{rem:periodogram-variance}
Lemma~\ref{lem:quadrature} bounds the \emph{bias}
\(\|\mathbb E[\widehat S^{\mathrm{xy}}_{\mathcal G_s}(u)]
- S^{\mathrm{xy}}(u)\|_{L^2}\) of the periodogram, not its
pointwise variance. From a single sample path the raw periodogram
has \(\Theta(1)\) variance per Fourier mode (chi-squared-2
distribution at each mode, cf.~\cite[Ch.~4]{Brillinger1981}); this
is the distributional content of Remark~\ref{rem:em-floor}. The
training loss (A4) therefore involves
\(\widehat S^{\mathrm{xy}}_{\mathcal G_s}(u_\theta)\), a
single-realisation quantity, while the bound
in~\eqref{eq:spectral-l2-bound} uses
\(\mathbb E[\widehat S^{\mathrm{xy}}_{\mathcal G_s}]\); the gap
between the two does not shrink as \(N_s \to \infty\) without
additional smoothing.

Lemma~\ref{lem:single-realisation} below closes this gap: replacing
the raw periodogram by its kernel-smoothed version
\(\widehat S^{\mathrm{xy}}_h\) with bandwidth \(h_N \to 0\),
\(N_s h_N^d \to \infty\), yields a consistent single-realisation
estimator at the minimax rate \(N_s^{-s_x/(2s_x+d)}\) under (A6).
In the practical regime of this paper (raw periodogram, single
realisation) we report against \(\gamma_{10}^2 \approx 0.063\) as
a single-mode reference scale (Remark~\ref{rem:em-floor}); the
surviving Gumbel-variance contribution and any residual law-distance
inflate the empirical \(\varepsilon\) of (A4) to the
\(5\)--\(7\times\) plateau seen in the experiments.
\end{remark}

\begin{lemma}[Single-realisation spectral consistency under mixing]
\label{lem:single-realisation}
Under (A1)--(A2) and (A6), let
\[
\widehat S^{\mathrm{xy}}_h(u)(k)
:= \int \kappa_{h_N}(k - k')\,
\widehat S^{\mathrm{xy}}_{\mathcal G_s}(u)(k')\,dk',
\]
where \(\kappa_{h_N}(k) = h_N^{-d}\kappa(k/h_N)\) is a symmetric
probability kernel with
\(\int |k|^{s_x}\kappa(k)\,dk < \infty\).
Then for any single realisation of \(u\),
\begin{equation}
\label{eq:sr-mse}
\bigl\|\widehat S^{\mathrm{xy}}_h(u) - S^{\mathrm{xy}}(u)
\bigr\|^2_{L^2(\widehat\Omega)}
\;\le\;
C_{\mathrm{bias}}\,h_N^{2s_x}
\;+\;
\frac{C_{\mathrm{var}}}{N_s\,h_N^{d}},
\end{equation}
where \(C_{\mathrm{bias}} = C_\kappa\|\mu^{\mathrm{xy}}\|^2_{H^{s_x}}\)
and \(C_{\mathrm{var}}\) depends on \(C_\alpha\) and \(\gamma\).
The minimax-optimal bandwidth \(h_N^* = N_s^{-1/(2s_x + d)}\)
gives the \(L^2\) rate
\[
\bigl\|\widehat S^{\mathrm{xy}}_{h^*}(u) - S^{\mathrm{xy}}(u)
\bigr\|_{L^2}
\;\le\;(C_{\mathrm{bias}}+C_{\mathrm{var}})^{1/2}\,
N_s^{-s_x/(2s_x+d)},
\]
which converges to zero and matches the minimax lower bound for
nonparametric spectral density estimation under \(\alpha\)-mixing
\cite{Dahlhaus1996}.
\end{lemma}
\begin{proof}
\emph{Bias.} Since
\(\mathbb E[\widehat S^{\mathrm{xy}}_{\mathcal G_s}(u)(k)] =
S^{\mathrm{xy}}(u)(k) + O(N_s^{-1})\) by
Lemma~\ref{lem:quadrature}, convolution with \(\kappa_{h_N}\)
gives
\(\mathbb E[\widehat S^{\mathrm{xy}}_h(u)(k)] =
(\kappa_{h_N} * S^{\mathrm{xy}})(k)\).
For \(S^{\mathrm{xy}} \in H^{s_x}\) and a kernel with vanishing
moments up to order \(s_x\), standard approximation theory gives
\(\|(\kappa_{h_N} * S^{\mathrm{xy}}) - S^{\mathrm{xy}}\|_{L^2}
\le C_\kappa \|\mu^{\mathrm{xy}}\|_{H^{s_x}} h_N^{s_x}\).

\emph{Variance.} Under (A6), Brillinger's cumulant spectral
bound~\cite[Thm.~4.3.2]{Brillinger1981} gives
\(|\mathrm{Cov}(\widehat S^{\mathrm{xy}}_{\mathcal G_s}(k_1),
\widehat S^{\mathrm{xy}}_{\mathcal G_s}(k_2))|
\le C N_s^{-1} \alpha(|k_1 - k_2|/\Delta k)\),
where \(\Delta k\) is the frequency resolution.
Convolving against \(\kappa_{h_N}^{\otimes 2}\) and using
\(\gamma > d + 1\) to ensure summability gives
\(\mathrm{Var}[\widehat S^{\mathrm{xy}}_h(u)(k)]
\le C_{\mathrm{var}}(N_s h_N^d)^{-1}\).
Integrating over \(k \in \widehat\Omega\) yields the variance term
in~\eqref{eq:sr-mse}.

\emph{Optimal bandwidth.} Minimising
\(C_{\mathrm{bias}} h_N^{2s_x} + C_{\mathrm{var}}(N_s h_N^d)^{-1}\)
over \(h_N\) gives \(h_N^* \propto N_s^{-1/(2s_x+d)}\) and the
stated rate. The minimax lower bound follows
from~\cite[Thm.~3.1]{Dahlhaus1996} applied to the stationary case.
\end{proof}

Combining Lemma~\ref{lem:log-linear} (applied to the smoothed
periodogram, per Remark~\ref{rem:why-smoothed}) with
Lemma~\ref{lem:single-realisation} (single-realisation consistency
of $\hat S^{\rm xy}_h$ at the minimax-optimal bandwidth
$h^\star = N_s^{-1/(2s_x+d)}$) via the triangle inequality, and
using $\|a+b\|^2 \le 2\|a\|^2 + 2\|b\|^2$:
\begin{equation}
\|S^{\mathrm{xy}}(u_\theta) - S^{\mathrm{xy}}(u_{\mathrm{FD}})
\|^2_{L^2(\widehat\Omega)}
\;\le\; 2(\ln 10)^2\,\mu_{\max}^2\;\varepsilon
\;+\; 4\,(C_{\mathrm{bias}}+C_{\mathrm{var}})\,
N_s^{-2 s_x / (2 s_x + d)},
\label{eq:spectral-l2-bound}
\end{equation}
which is the rigorous single-realisation bound; the $N_s^{-2s_x/(2s_x+d)}$
exponent gives $\beta_{\rm sr} = s_x/(2(2s_x+d))$ in
Proposition~\ref{thm:main}. An analogous bound holds for the
temporal direction.

\begin{remark}[Bias-only rate]
\label{rem:bias-only-rate}
If the realisation variance of the periodogram is ignored
(equivalently, the chain replaces
Lemma~\ref{lem:single-realisation} by Lemma~\ref{lem:quadrature}'s
bias bound), one obtains the faster bias-only rate
$2(\ln 10)^2 \mu_{\max}^2 \varepsilon + 2 C_{\rm quad,x}(N_s^x)^{-2 s_x/d}$
matching $\beta_{\rm bias} = \tfrac12\min(s_x/d,\,s_t)$ (which
collapses to $s_x/(2d)$ when $s_t \ge s_x/d$, the regime of this
paper). This rate corresponds to
training against an ensemble-averaged spectral target rather than a
single realisation, and is the rate observed empirically in
Section~\ref{subsec:grid-conv} of the main text where the target is
ensemble-averaged. The rigorous single-realisation rate
$\beta_{\rm sr} = 1/8$ is slower but applies to any individual
training run; the bias-only rate $\beta_{\rm bias} = 1/4$ is faster
but is a heuristic statement about the bias of the spectral target,
not about training-error convergence per realisation.
\end{remark}

\begin{lemma}[Plancherel--Gelbrich bound]
\label{lem:gelbrich}
Let \(u_1, u_2\) be centred Gaussian random fields on
\(\Omega \times [0,T]\) satisfying (A1)--(A3). Then
\[
W_2^2\bigl(\mathcal L(u_1),\,\mathcal L(u_2)\bigr)
\;\le\; \frac{1}{4\mu_{\min}}\;
\|\mu_1 - \mu_2\|^2_{L^2(\widehat\Omega \times [0,T])}.
\]
\end{lemma}
\begin{proof}
For centred Gaussian measures on a separable Hilbert space with
covariance operators \(C_1, C_2\), the Gelbrich--Givens--Shortt
formula~\cite{Gelbrich1990,GivensShortt1984} gives
\(W_2^2 = \mathrm{tr}\bigl(C_1 + C_2
- 2(C_1^{1/2} C_2 C_1^{1/2})^{1/2}\bigr)\).
By (A1), both fields are stationary on the periodised domain;
hence \(C_1, C_2\) are simultaneously diagonalised by the Fourier
basis with eigenvalues \(\mu_1(k,\omega), \mu_2(k,\omega)\). The
formula reduces to
\(W_2^2 = \sum_{(k,\omega)}
(\sqrt{\mu_1(k,\omega)} - \sqrt{\mu_2(k,\omega)})^2\).
From the identity
\((\sqrt a - \sqrt b)^2 = (a-b)^2/(\sqrt a + \sqrt b)^2\)
and \((\sqrt a + \sqrt b)^2 \ge 4\min(a,b)\), together with (A3):
\(W_2^2 \le \frac{1}{4\mu_{\min}}
\sum_{(k,\omega)} (\mu_1(k,\omega) - \mu_2(k,\omega))^2\).
Under separability, the tensor-product decomposition
\(\mu_1 - \mu_2 = (\mu_1^{\mathrm{xy}} - \mu_2^{\mathrm{xy}})
\otimes \mu_1^t + \mu_2^{\mathrm{xy}} \otimes
(\mu_1^t - \mu_2^t)\) and \(\|a+b\|^2 \le 2\|a\|^2 + 2\|b\|^2\)
give the factored form used in the proof of
Proposition~\ref{thm:main}.
\end{proof}

\begin{lemma}[Mishra--Molinaro, adapted]
\label{lem:mm}
Assume the SPDE residual operator \(\mathcal R\) is well posed in the
sense of \cite[Thm.~3.4]{MishraMolinaro2022}, with stability constant
\(C_{\mathrm{stab}}\). Then under \eqref{eq:tolerance},
\[
\mathbb E\,\|u_\theta - u_{\mathrm{FD}}\|^2_{L^2(\Omega \times [0,T])}
\;\le\; C_{\mathrm{MM}}\,\delta,
\quad
C_{\mathrm{MM}} = C_{\mathrm{stab}}
\bigl(1 + \lambda_{\mathrm{ic}}^{-1}
+ \lambda_{\mathrm{bc}}^{-1}\bigr).
\]
\end{lemma}
\begin{proof}
The driver fields \(\xi, J_{\mathrm{VG}}\) are identical between
\(u_\theta\) and \(u_{\mathrm{FD}}\) (spline-interpolated from the
same FD sample path), so the residual operator reduces to a
deterministic PDE error analysis in the realisation
\(\omega \in \mathcal P\). Apply~\cite[Thm.~3.4]{MishraMolinaro2022}
pathwise and take expectations. Since
\(W_2^2(\mathcal L(u_1), \mathcal L(u_2))
\le \mathbb E\,\|u_1 - u_2\|^2\) for the trivial coupling, this
contributes \(\sqrt{C_{\mathrm{MM}}\,\delta}\) to the \(W_2\) bound.

\emph{Remark on the spline-smoothing error.} The spline
interpolation replaces jump discontinuities in
\(\xi, J_{\mathrm{VG}}\) with smooth approximants, introducing a
smoothing error not explicitly bounded in \(C_{\mathrm{MM}}\). For
Brownian forcing the interpolation error vanishes as $N_s \to
\infty$ at the standard $O(N_s^{-1/2})$ rate and is absorbed into
$\delta$ without loss. For Variance-Gamma forcing
(unbounded-variation paths) the jump-component interpolation error
is $\Theta(1)$ as $N_s \to \infty$: the spline misses the jumps
regardless of grid refinement. The bound
$\mathbb E\|u_\theta - u_{\mathrm{FD}}\|^2 \le C_{\rm MM}\delta$
therefore degrades to
$\le C_{\rm MM}\delta + \mathrm{(jump\;floor)}$ for VG, where the
jump floor is not controlled by training. This is the principal
analytical obstacle to a fully rigorous L\'evy extension and
motivates the partial-result framing of
Section~\ref{subsec:levy-ext} in the main text. The rigorous main
theorem (Proposition~\ref{thm:main}) assumes (A5) Gaussianity, so
this remark applies only to the L\'evy extensions
(Propositions~\ref{thm:levy-ext}, \ref{prop:levy-ito}).
\end{proof}

\begin{proof}[Proof of Proposition~\ref{thm:main}]
The bound is obtained as the looser of two independent upper bounds
on \(W_2(\mathcal L(u_\theta),\mathcal L(u_{\mathrm{FD}}))\), combined
via \(\min(a,b)\le a+b\).

\smallskip
\emph{Spectral path (uses (A5)).} Apply the triangle inequality on
\(W_2\) through the intermediate centred Gaussian measure
\(\mathcal N(0,\,\mathcal C_{u_{\mathrm{FD}}})\) with covariance
operator \(\mathcal C_{u_{\mathrm{FD}}}\) equal to that of
\(u_{\mathrm{FD}}\):
\[
W_2\bigl(\mathcal L(u_\theta),\,\mathcal L(u_{\mathrm{FD}})\bigr)
\;\le\;
W_2\bigl(\mathcal L(u_\theta),\,
\mathcal N(0,\,\mathcal C_{u_{\mathrm{FD}}})\bigr)
\,+\,
W_2\bigl(\mathcal N(0,\,\mathcal C_{u_{\mathrm{FD}}}),\,
\mathcal L(u_{\mathrm{FD}})\bigr).
\]
Under (A5) both fields are centred Gaussian, so
\(\mathcal L(u_{\mathrm{FD}}) = \mathcal N(0,\,\mathcal C_{u_{\mathrm{FD}}})\)
and the second term vanishes;
\(\mathcal L(u_\theta) = \mathcal N(0,\,\mathcal C_{u_\theta})\) and
Lemma~\ref{lem:gelbrich} bounds the first term by
\(\tfrac{1}{2\sqrt{\mu_{\min}}}\,\|\mu_{u_\theta}
- \mu_{u_{\mathrm{FD}}}\|_{L^2}\). Using the separable factoring
\(\mu_1 - \mu_2 = (\mu^{\mathrm{xy}}_1 - \mu^{\mathrm{xy}}_2)
\otimes \mu^t_1
+ \mu^{\mathrm{xy}}_2 \otimes (\mu^t_1 - \mu^t_2)\) with
\(\|a+b\|^2 \le 2\|a\|^2 + 2\|b\|^2\), and bounding the trained
model's marginal norm $\|\mu^t_1\|_{L^2} \le M_t + \|\mu^t_1 -
\mu^t_2\|_{L^2}$ (triangle), the second summand is itself at most
$2(\ln 10)\mu_{\max}\sqrt\varepsilon$ by
Lemma~\ref{lem:log-linear}, which is a lower-order contribution
absorbed into the $\sqrt\varepsilon$ term; treating it as such, we
may use $M_t$ uniformly in the factoring. Substituting
\eqref{eq:spectral-l2-bound} for each spectral $L^2$ difference at
$h^\star$:
\[
W_2 \;\le\;
\frac{1}{\sqrt{2\mu_{\min}}}\Bigl[
2(\ln 10)\,\mu_{\max}\sqrt\varepsilon\,(M_t + M_{xy})
+ 2\sqrt{(C_{\rm bias}+C_{\rm var})}\,N_s^{-s_x/(2 s_x + d)}\,M_t
+ 2\sqrt{(C_{\rm bias,t}+C_{\rm var,t})}\,N_s^{-s_t/(2 s_t + 1)}\,M_{xy}
\Bigr],
\]
where \(M_t := \|\mu^t_{\mathrm{FD}}\|_{L^2}\),
\(M_{xy} := \|\mu^{\mathrm{xy}}_{\mathrm{FD}}\|_{L^2}\), and the
spatial and temporal bandwidths $h^\star_x = N_s^{-1/(2s_x+d)}$,
$h^\star_t = N_s^{-1/(2s_t+1)}$ are set to their respective
minimax-optimal values. Absorbing constants and using
\(N_s^{-\beta_{\rm sr}} \ge N_s^{-s_x/(2 s_x + d)}
+ N_s^{-s_t/(2 s_t + 1)}\) up to a constant
(definition of $\beta_{\rm sr} = s_x/(2(2s_x+d))$ at $s_t \ge s_x$),
\[
W_2\bigl(\mathcal L(u_\theta),\,\mathcal L(u_{\mathrm{FD}})\bigr)
\;\le\; C_1\sqrt\varepsilon \,+\, C_2\,N_s^{-\beta_{\rm sr}}.
\tag{$\star$}\label{eq:spectral-path}
\]

\smallskip
\emph{Pathwise path (does not use (A5)).} Lemma~\ref{lem:mm} applied
with the trivial coupling gives
\[
W_2^2\bigl(\mathcal L(u_\theta),\,\mathcal L(u_{\mathrm{FD}})\bigr)
\;\le\; \mathbb E\,\|u_\theta - u_{\mathrm{FD}}\|^2
\;\le\; C_{\mathrm{MM}}\,\delta,
\quad\text{hence}\quad
W_2 \;\le\; C_3\sqrt\delta,
\tag{$\dagger$}\label{eq:pde-path}
\]
with no Gaussianity assumption.

\smallskip
\emph{Combined bound.} Since \(\min(a,b) \le a + b\) for
\(a,b\ge 0\), \eqref{eq:spectral-path} and \eqref{eq:pde-path} together
yield
\[
W_2\bigl(\mathcal L(u_\theta),\,\mathcal L(u_{\mathrm{FD}})\bigr)
\;\le\; C_1\sqrt\varepsilon \,+\, C_2\,N_s^{-\beta_{\rm sr}}
\,+\, C_3\sqrt\delta,
\]
which is \eqref{eq:main-bound}. The sum form is loose; in either
asymptotic regime (spectrally-rich training with vanishing
\(\varepsilon\), or PDE-tight training with vanishing \(\delta\))
the corresponding single-path bound is the operative one. The
spectral path requires (A5); the PDE path does not. In particular,
dropping the spectral loss from training (so the bound on
\(\varepsilon\) becomes vacuous) reduces the result to
\eqref{eq:pde-path}, the Mishra--Molinaro estimate of
\cite{MishraMolinaro2022}.
\end{proof}

\subsubsection{Corollaries}

\begin{proposition}[Variance-gap rate]
\label{cor:vr}
Under (A1)--(A4),
\[
\bigl|\mathrm{Var}(u_\theta)/\mathrm{Var}(u_{\mathrm{FD}}) - 1\bigr|
\;\le\;
\frac{2\sqrt{2}}{\sqrt{\mathrm{Var}(u_{\mathrm{FD}})}}
\bigl(\sqrt{\mathrm{Var}(u_\theta)}
+ \sqrt{\mathrm{Var}(u_{\mathrm{FD}})}\bigr)
\bigl(C_1\sqrt\varepsilon + C_2 N_s^{-\beta_{\rm sr}}
+ C_3\sqrt\delta\bigr).
\]
\end{proposition}
\begin{proof}
For centred fields, \(|\mathrm{Var}(u_1) - \mathrm{Var}(u_2)|
= |\mathbb E[\|u_1\|^2] - \mathbb E[\|u_2\|^2]|\).
Under the optimal coupling \((X, Y)\) with \(X \sim \mathcal L(u_\theta)\),
\(Y \sim \mathcal L(u_{\mathrm{FD}})\):
\[
|\mathbb E[\|X\|^2] - \mathbb E[\|Y\|^2]|
\le \mathbb E\bigl[({\|X\|+\|Y\|})\,\|X-Y\|\bigr]
\le \sqrt{2(\mathrm{Var}(u_\theta)+\mathrm{Var}(u_{\mathrm{FD}}))}
\cdot W_2
\le \sqrt{2}\,(\sqrt{\mathrm{Var}(u_\theta)}
+\sqrt{\mathrm{Var}(u_{\mathrm{FD}})})\cdot W_2,
\]
using Cauchy--Schwarz and \((a+b)^2\le 2(a^2+b^2)\).
Divide by \(\mathrm{Var}(u_{\mathrm{FD}})\) and substitute
\eqref{eq:main-bound}.
\end{proof}

\begin{remark}[Reference scale \(\gamma_{10}^2\) for the
raw-periodogram loss]
\label{rem:em-floor}
The raw log-periodogram is asymptotically Gumbel with mean
\(\log\mu(k) - \gamma\) and per-mode variance
\(\pi^2/6\)~\cite[Ch.~5]{Brillinger1981}. In our \(\log_{10}\)
formulation the shift is \(\gamma_{10} := \gamma/\ln 10 \approx 0.251\);
it is identical for predictor and reference and therefore cancels
exactly in the prediction--reference difference, contributing
nothing to \(\varepsilon\) in (A4) in expectation. The non-vanishing
residual is the per-mode Gumbel variance
\(\pi^2/(6\ln^2\!10) \approx 0.310\), damped by axis-averaging in
\(\widehat S_{xy}, \widehat S_t\) but non-zero at any fixed \(N_s\).
We report empirical losses against the single-mode reference scale
\(\gamma_{10}^2 \approx 0.063\) --- the squared mean bias of one
log-periodogram against the true PSD --- and observe per-channel
log-PSD MSE plateauing near \(0.3\)--\(0.4\) (a \(5\)--\(7\times\)
multiplier on \(\gamma_{10}^2\)), the multiplier absorbing both the
surviving variance contribution and the residual law-distance. The
corresponding reference \(W_2\) contribution is \(C_1\gamma_{10}\).
The kernel-smoothed estimator of
Lemma~\ref{lem:single-realisation} drives the variance residual to
zero at rate \(N_s^{-2s_x/(2s_x+d)}\), so no \(\gamma_{10}\)
benchmark is needed in the smoothed proof framework.
\end{remark}

\subsubsection{Extension to L\'evy-driven fields}

\begin{proposition}[Second-order convergence for stationary L\'evy fields]
\label{thm:levy-ext}
Under (A1)--(A4) but with \(u_\theta, u_{\mathrm{FD}}\) not
necessarily Gaussian, denote by \(\Pi_2\mathcal L(u)\) the Gaussian
projection (the centred Gaussian measure with the same covariance).
Then
\[
W_2\bigl(\Pi_2\mathcal L(u_\theta),\,
\Pi_2\mathcal L(u_{\mathrm{FD}})\bigr)
\;\le\; C_1\sqrt\varepsilon + C_2 N_s^{-\beta_{\rm sr}} + C_3\sqrt\delta.
\]
\end{proposition}
\begin{proof}
Identical to the proof of Proposition~\ref{thm:main}, replacing
\(\mathcal L(u)\) by \(\Pi_2\mathcal L(u)\) throughout. The spectral
measure \(\mu_u\) depends only on the covariance, hence is invariant
under the projection.
\end{proof}

The gap between \(\mathcal L(u)\) and \(\Pi_2\mathcal L(u)\) is
determined by the cumulants \(\kappa_k(u)\) for \(k \ge 3\). For
Variance-Gamma drivers \(\kappa_3, \kappa_4\) are non-zero, so
Proposition~\ref{thm:levy-ext} does \emph{not} imply convergence in
\(W_2\) on the full law. Two routes close this gap:
\begin{enumerate}
\item a cumulant-matching loss
$\mathcal{L}_{\kappa_4}(u_\theta):=(\kappa_4(u_\theta)-\kappa_4(u_{\mathrm{FD}}))^2$
(Section~\ref{subsec:cumulant-loss}), which empirically raises the
predicted kurtosis from $4.2$ to $25.1$; or
\item an $L^2$ bound on the pure-jump component $u^J$ via the
L\'evy--It\^o decomposition (Proposition~\ref{prop:levy-ito} below),
giving an explicit $O(\sqrt\nu)$ correction to~\eqref{thm:levy-ext}.
\end{enumerate}

\begin{proposition}[Full-law convergence via L\'evy--It\^o coupling]
\label{prop:levy-ito}
Under (A1)--(A4) and (A7),
\begin{equation}
\label{eq:levy-ito-bound}
W_2\!\left(\mathcal{L}(u_\theta),\,\mathcal{L}(u_{\mathrm{FD}})\right)
\;\le\;C_1\sqrt{\varepsilon}+C_2 N_s^{-\beta_{\rm sr}}+C_3\sqrt{\delta}
+4\sqrt{J}.
\end{equation}
In the Gaussian limit $J\to 0$ (equivalently $\nu\to 0$ for VG
drivers), this reduces to Proposition~\ref{thm:main}.
\end{proposition}

\begin{proof}
Apply the $W_2$ triangle inequality:
\begin{align}
\label{eq:levy-ito-triangle}
W_2(\mathcal{L}(u_\theta),\mathcal{L}(u_{\mathrm{FD}}))
&\le
W_2\!\left(\mathcal{L}(u_\theta),\Pi_2\mathcal{L}(u_\theta)\right)
+W_2\!\left(\Pi_2\mathcal{L}(u_\theta),\Pi_2\mathcal{L}(u_{\mathrm{FD}})\right) \notag\\
&\quad+W_2\!\left(\Pi_2\mathcal{L}(u_{\mathrm{FD}}),\mathcal{L}(u_{\mathrm{FD}})\right).
\end{align}
The middle term is bounded by
$C_1\sqrt\varepsilon+C_2 N_s^{-\beta_{\rm sr}}+C_3\sqrt\delta$
by Proposition~\ref{thm:levy-ext}.
Write $u=u^G+u^J$ as in (A7), with
$u^G\sim\mathcal{N}(0,C_{u^G})$ and
$\Pi_2\mathcal{L}(u)=\mathcal{N}(0,C_u)=\mathcal{N}(0,C_{u^G}+C_{u^J})$.
Couple $(\mathcal{L}(u),\mathcal{N}(0,C_u))$ by
$X:=u^G+u^J$ and $Y:=u^G+G_2$
where $G_2\sim\mathcal{N}(0,C_{u^J})$ is independent of $u$.
Then $Y\sim\mathcal{N}(0,C_u)$ and
\[
\mathbb{E}\|X-Y\|^2_{L^2}
=\mathbb{E}\|u^J-G_2\|^2_{L^2}
\le 2\mathbb{E}\|u^J\|^2_{L^2}+2\,\mathrm{tr}(C_{u^J})
= 4\mathbb{E}\|u^J\|^2_{L^2}\le 4J,
\]
giving $W_2(\mathcal{L}(u),\Pi_2\mathcal{L}(u))\le 2\sqrt{J}$ for
each of $u_\theta$ and $u_{\mathrm{FD}}$.
Inserting into~\eqref{eq:levy-ito-triangle} yields~\eqref{eq:levy-ito-bound}.
$\hfill\square$
\end{proof}

\begin{remark}[VG quantification of $J$]
\label{rem:vg-J}
For the Variance-Gamma driver with parameters $(\sigma,\nu,\theta)$,
the pure-jump component satisfies
$\mathbb{E}\|\xi_{\mathrm{VG}}^J\|^2_{L^2([0,T])}=\nu(\sigma^2+\nu\theta^2)T(1+O(\nu))$.
The SPDE solution inherits this via the Green's function, giving
$J=O(\nu)$ and a correction $4\sqrt{J}=O(\sqrt\nu)$ in
\eqref{eq:levy-ito-bound} that vanishes in the Gaussian limit.
The relationship $\kappa_4=3\nu$ links the correction directly to the
fourth cumulant: $4\sqrt{J}=O(\kappa_4^{1/2}/3^{1/2})$.
\end{remark}

\begin{proposition}[Berry--Esseen tightness of spectral matching]
\label{cor:berry-esseen}
Under (A1)--(A4) and (A8) with $m=3$,
\begin{equation}
\label{eq:be-bound}
W_1\!\left(\mathcal{L}(u(x,t)),\;\mathcal{N}(0,\sigma_u^2(x,t))\right)
\;\le\;
C_{\mathrm{BE}}\,\frac{|\kappa_3(u(x,t))|}{\sigma_u^2(x,t)},
\end{equation}
where $C_{\mathrm{BE}}$ is a universal constant from Stein's
method~\cite{ChenGoldsteinShao2011}. For Variance-Gamma drivers with
parameters $(\sigma_{\mathrm{vg}},\nu,\theta)$, the exact third cumulant
per unit volume is $\kappa_3(\xi)=\nu\theta(3\sigma_{\mathrm{vg}}^2+2\nu\theta^2)$,
giving
\begin{equation}
\label{eq:kappa3-vg}
|\kappa_3(u(x,t))|\;\le\;
C_G^3\,\nu\,|\theta|\,(3\sigma_{\mathrm{vg}}^2+2\nu\theta^2),\quad
C_G:=\|G_{\mathrm{SPDE}}(x-\cdot,t-\cdot)\|_{L^3},
\end{equation}
and hence, for $2\nu\theta^2\ll 3\sigma_{\mathrm{vg}}^2$,
\begin{equation}
\label{eq:be-vg}
W_1\!\left(\mathcal{L}(u(x,t)),\;\mathcal{N}(0,\sigma_u^2)\right)
\;\lesssim\;
\frac{3\,C_{\mathrm{BE}}\,C_G^3\,\nu\,\sigma_{\mathrm{vg}}^2\,|\theta|}{\sigma_u^2}.
\end{equation}
\paragraph{Tightness condition.} Spectral matching controls $\sigma_u$
(second moment); the residual non-Gaussianity is at most
$\varepsilon_{\mathrm{gap}}$ in $W_1$ whenever
$|\theta|\le\theta_{\mathrm{cut}}:=
\varepsilon_{\mathrm{gap}}\,\sigma_u^2\,/\,
(3\,C_{\mathrm{BE}}\,C_G^3\,\nu\,\sigma_{\mathrm{vg}}^2)$.
\end{proposition}

\begin{proof}
\eqref{eq:be-bound} is the Stein's-method bound
$W_1(\mathcal{L}(X),\mathcal{N}(0,\sigma^2))\le C_S\,E|X|^3/\sigma^2$
for a centred $X$~\cite[Thm.~3.1]{ChenGoldsteinShao2011},
combined with $|\kappa_3|=|E[X^3]|\le E|X|^3$.
For~\eqref{eq:kappa3-vg}: by stochastic Fubini,
$\kappa_3(u)=\int G^3\kappa_3(\xi)\,dy\,ds$
(multilinearity of cumulants under independence of Lévy increments).
Expanding $\ln\phi(u)=-(1/\nu)\ln(1-u\theta\nu-u^2\sigma_{\mathrm{vg}}^2\nu/2)$
to third order gives $\kappa_3(\xi)=\nu\theta(3\sigma_{\mathrm{vg}}^2+2\nu\theta^2)$~\cite{ContTankov2003,MadanCarrChang1998}.
H\"older's inequality then gives~\eqref{eq:kappa3-vg}; the dominant small-$|\theta|$
term yields~\eqref{eq:be-vg}.
$\hfill\square$
\end{proof}

\begin{remark}[Numerical tightness]
\label{rem:be-numerical}
At $\sigma_{\mathrm{vg}}=0.2$, $\nu=0.5$, $\sigma_u\approx 0.3$,
$C_{\mathrm{BE}}\approx 1$, $C_G=O(1)$:
$\theta_{\mathrm{cut}}(0.05)
\approx 0.09/(0.06\,C_G^3)
\approx 0.075/C_G^3$.
Near-symmetric VG with $|\theta|\lesssim 0.07$ gives $5\%$ tightness.
The paper's $\theta=-0.1$ produces a $W_1$ residual $\approx 0.07/C_G^3$;
the kurtosis gap beyond this is targeted by the cumulant-matching loss
of Section~\ref{subsec:cumulant-loss}.
\end{remark}

\begin{remark}[Single-realisation rate vs bias-only rate]
\label{cor:smoothed-loss}
Proposition~\ref{thm:main}'s single-realisation rate
$\beta_{\rm sr} = s_x/(2(2 s_x + d))$ is obtained by combining
Lemma~\ref{lem:single-realisation} (smoothed-periodogram
consistency) with the spectral $L^2$ chain
\eqref{eq:spectral-l2-bound}. The faster bias-only rate
$\beta_{\rm bias} = \tfrac12\min(s_x/d,\,s_t)$ from
Remark~\ref{rem:bias-only-rate}
corresponds to replacing Lemma~\ref{lem:single-realisation} by
Lemma~\ref{lem:quadrature}'s expected-periodogram bound, which is
appropriate when the spectral target is averaged across
realisations. For $s_x = s_t = 1, d = 2$: $\beta_{\rm sr} = 1/8$
vs $\beta_{\rm bias} = 1/4$. The slower exponent is the price of
single-realisation validity. The full proof of
Proposition~\ref{thm:main} uses the single-realisation rate
throughout; the bias-only rate is kept only as connection with the
empirical grid-convergence study.
\end{remark}

\begin{proof}[Sketch of the relation between the two rates]
Replace Lemma~\ref{lem:quadrature} in the proof of
Proposition~\ref{thm:main} with Lemma~\ref{lem:single-realisation}
at \(h_N = h_N^*\). The \(L^2\) spectral error for a single
realisation is
\((C_{\mathrm{bias}}+C_{\mathrm{var}})^{1/2} N_s^{-s_x/(2s_x+d)}\).
Inserting into the Gelbrich step (Lemma~\ref{lem:gelbrich}), which
bounds \(W_2 \le \tfrac{1}{2\sqrt{\mu_{\min}}}
\|\mu_1 - \mu_2\|_{L^2}\), gives the factor
\(N_s^{-s_x/(2(2s_x+d))} = N_s^{-\beta_{\mathrm{sr}}}\).
\end{proof}

\subsubsection{Extension to non-separable covariance}
\label{subsec:nonsep-extension}

The separability assumption (A1) is load-bearing for two steps of the
proof: the tensor-product factoring of \(\mu_1 - \mu_2\) in
Lemma~\ref{lem:gelbrich}, and the design of the spectral loss, which
matches \(\mu^{\mathrm{xy}}\) and \(\mu^t\) separately rather than
the joint measure \(\mu(\omega, k)\). Both restrictions can be lifted
by working with the \emph{joint} \((d+1)\)-dimensional spectral
density, at the cost of a weaker quadrature rate.

\paragraph{Relaxed separability assumption.}
Replace (A1) by:
\begin{description}
\item[(A1$'$) Stationarity without separability.]
Both \(u_\theta\) and \(u_{\mathrm{FD}}\) are centred, weakly
stationary with absolutely continuous joint spectral density
\(S^{\mathrm{joint}}(\omega, k) \in H^s(\mathbb R^{d+1})\),
\(s = \min(s_x, s_t) > (d+1)/2\). No tensor-product factorisation
is assumed.
\end{description}

The joint spectral loss replaces
\(\mathcal L_{\mathrm{xy}} + \mathcal L_t\) in (A4) by
\begin{equation}
\label{eq:joint-tolerance}
\mathcal L_{\mathrm{joint}}(u_\theta) \le \varepsilon', \quad
\mathcal L_{\mathrm{joint}}(u) =
\overline{\Bigl(
\log_{10}\widehat S^{\mathrm{joint}}_{\mathcal G_s}(u)
- \log_{10}\widehat S^{\mathrm{joint}}_{\mathcal G_s}(u_{\mathrm{FD}})
\Bigr)^2},
\tag{A4$'$}
\end{equation}
where
\(\widehat S^{\mathrm{joint}}_{\mathcal G_s}(u)(\omega, k) =
|\mathcal G_s|^{-1}
|\mathcal F_{x,t}[u\!\upharpoonright_{\mathcal G_s}](\omega,k)|^2\)
is the \((d+1)\)-dimensional periodogram on \(\mathcal G_s\).

\begin{lemma}[Joint spectral quadrature error]
\label{lem:joint-quadrature}
Under (A1$'$), for every stationary \(u\) with
\(S^{\mathrm{joint}} \in H^s(\mathbb R^{d+1})\), \(s > (d+1)/2\),
\[
\bigl\|\mathbb E\bigl[\widehat S^{\mathrm{joint}}_{\mathcal G_s}(u)\bigr]
- S^{\mathrm{joint}}(u)\bigr\|^2_{L^2(\mathbb R^{d+1})}
\;\le\; C_{\mathrm{quad}}\,(N_s^{\mathrm{tot}})^{-2s/(d+1)},
\]
where \(N_s^{\mathrm{tot}} = N_s^x N_s^T\) is the total node count
and \(C_{\mathrm{quad}} = C\,\|S^{\mathrm{joint}}\|^2_{H^s}\) with
\(C\) depending only on \(\mathrm{diam}(\Omega)\), \(T\), and \(d\).
\end{lemma}
\begin{proof}
The expected \((d+1)\)-dimensional periodogram is a Dirichlet-kernel
smoothing of \(S^{\mathrm{joint}}\) in \(\mathbb R^{d+1}\), inheriting
its \(H^s\) regularity. The \(L^2\) trapezoidal-rule error rate for
\(H^s\) functions on a \((d+1)\)-dimensional torus
(cf.~\cite{Brillinger1981,Trefethen2000}) with spacing
\(h = (N_s^{\mathrm{tot}})^{-1/(d+1)}\) gives bias
\(O(h^s) = O((N_s^{\mathrm{tot}})^{-s/(d+1)})\) in \(L^2\); squaring
yields the stated bound.
\end{proof}

\begin{proposition}[Convergence rate under non-separable covariance]
\label{thm:nonsep}
Under (A1$'$), (A3), (A4$'$), (A5) with \(s = \min(s_x, s_t)\),
\begin{equation}
\label{eq:nonsep-bound}
W_2\bigl(\mathcal L(u_\theta),\,\mathcal L(u_{\mathrm{FD}})\bigr)
\;\le\; C_1'\sqrt{\varepsilon'}
\;+\; C_2'\,(N_s^{\mathrm{tot}})^{-\beta'}
\;+\; C_3\sqrt{\delta},
\end{equation}
where
\[
\beta' = \frac{s}{2(d+1)},
\]
and \(C_1' = \mu_{\max}\,\ln 10 / \sqrt{2\,\mu_{\min}}\),
\(C_2' = \sqrt{C_{\mathrm{quad}}/\mu_{\min}}\),
\(C_3 = \sqrt{C_{\mathrm{MM}}}\).
Compared to the separable Proposition~\ref{thm:main} the joint-density
constants have no tensor-product factors \(M_{xy},M_t\); the
\(\sqrt{2}\) in the denominator is from the
\((\sqrt{a}+\sqrt{b})^2 \ge 4\min(a,b)\) step in
Lemma~\ref{lem:gelbrich}.
\end{proposition}
\begin{proof}
\emph{Spectral path (uses (A5)).} Apply the triangle inequality
through \(\mathcal N(0, \mathcal C_{u_{\mathrm{FD}}})\) as in
Proposition~\ref{thm:main}; under (A5) the second term vanishes.
Lemma~\ref{lem:gelbrich} gives
\(W_2 \le \tfrac{1}{2\sqrt{\mu_{\min}}}
\|\mu_{u_\theta} - \mu_{u_{\mathrm{FD}}}\|_{L^2(\mathbb R^{d+1})}\).
Since (A1$'$) does not assume separability there is no
tensor-product factoring; apply the triangle inequality directly:
\begin{align*}
\|\mu_1 - \mu_2\|^2_{L^2}
\;\le\;&
2\,\bigl\|\mu_1 -
\mathbb E[\widehat S^{\mathrm{joint}}_{\mathcal G_s}(u_\theta)]\bigr\|^2_{L^2}
+ 2\,\bigl\|\mu_2 -
\mathbb E[\widehat S^{\mathrm{joint}}_{\mathcal G_s}(u_{\mathrm{FD}})]\bigr\|^2_{L^2}
\\
&+ 2\,\bigl\|\widehat S^{\mathrm{joint}}_{\mathcal G_s}(u_\theta)
- \widehat S^{\mathrm{joint}}_{\mathcal G_s}(u_{\mathrm{FD}})\bigr\|^2_{L^2}.
\end{align*}
The first two terms are bounded by
\(2C_{\mathrm{quad}}(N_s^{\mathrm{tot}})^{-2s/(d+1)}\)
(Lemma~\ref{lem:joint-quadrature}); the third by
\(2(\ln 10)^2\mu_{\max}^2\varepsilon'\) via
Lemma~\ref{lem:log-linear} and~\eqref{eq:joint-tolerance}. Taking
square roots:
\[
W_2 \;\le\; C_1'\sqrt{\varepsilon'} + C_2'(N_s^{\mathrm{tot}})^{-\beta'}.
\tag{$\star'$}\label{eq:spectral-path-nonsep}
\]
\emph{Pathwise path.} Lemma~\ref{lem:mm} is unchanged:
\(W_2 \le C_3\sqrt\delta\).

\emph{Combined.} \(\min(a,b) \le a + b\) yields \eqref{eq:nonsep-bound}.
\end{proof}

\paragraph{Rate comparison.}
The exponent \(\beta' = s/(2(d+1))\) is always weaker than the
separable bias-only rate
\(\beta_{\rm bias} = \tfrac12\min(s_x/d,\,s_t)\): separability permits
estimating two lower-dimensional marginals (\(\mu^{\mathrm{xy}}\) in
\(\mathbb R^d\), \(\mu^t\) in \(\mathbb R\)) rather than a single
density in \(\mathbb R^{d+1}\), exploiting additional structure. Both
\(\beta'\) and \(\beta_{\rm bias}\) are ensemble-target exponents; the
single-realisation analogue of \(\beta'\) is not pursued here.
Concrete values:
\begin{center}
\begin{tabular}{lcc}
\toprule
Setting & Separable \(\beta_{\rm bias}\) & Non-sep.\ \(\beta'\) \\
\midrule
\(s=1,\,d=2\) & \(\tfrac14\) & \(\tfrac16\) \\
\(s=2,\,d=2\) & \(\tfrac12\) & \(\tfrac13\) \\
\(s=1,\,d=3\) & \(\tfrac16\) & \(\tfrac18\) \\
\bottomrule
\end{tabular}
\end{center}
Despite the weaker rate, Proposition~\ref{thm:nonsep} is the
appropriate result when the SPDE solution's covariance does not
factorise, the generic case for L\'evy-driven equations with
space-time-coupled forcing.

\subsubsection{Discussion of the rate}

\paragraph{Optimality.}
The single-realisation exponent
\(\beta_{\rm sr} = s_x/(2(2s_x+d))\) is the minimax-optimal rate for
non-parametric spectral density estimation under
$\alpha$-mixing~\cite{Dahlhaus1996} and is sharp without
strengthening (A2) to analyticity. The bias-only exponent
\(\beta_{\rm bias} = \tfrac12\min(s_x/d, s_t)\) coincides with the
standard FFT quadrature rate for $H^s$ functions on a torus and is
also sharp for the bias-only target. Under analytic $\mu$ the
rates improve to spectral $\exp(-c\,N_s^{1/d})$, but this is not
the relevant regime for L\'evy-driven fields.

\paragraph{Floor regime.}
For finite \(\varepsilon, \delta\) the \(\sqrt\varepsilon\) term
dominates once \(N_s \gg \varepsilon^{-1/(2\beta_{\rm sr})}\). At the
paper's operating point
\(N_s = 24 \times 24 \times 96 \approx 5.5 \times 10^4\) with
\(s_x = s_t = 1\), \(d = 2\) and hence \(\beta_{\rm sr} = 1/8\), the
quadrature factor is \(N_s^{-\beta_{\rm sr}} \approx 0.25\) (versus the
bias-only $N_s^{-\beta_{\rm bias}} \approx 0.065$ that would apply
under ensemble averaging of the spectral target). For typical
\(\varepsilon \approx 0.3\) the training-error term contributes
\(\sim \mu_{\max}\sqrt{0.3}/\sqrt{\mu_{\min}}\), which is the
dominant source of \(W_2\) error in this regime; the quadrature
term is smaller by a factor depending on the unresolved
\(\mu_{\max}/\sqrt{\mu_{\min}}\) ratio but does not vanish.

\paragraph{Connection to numerical results.}
For Burgers-VG (\(s_x = s_t = 1\), \(d = 2\),
\(N_s \approx 5.5 \times 10^4\)),
Proposition~\ref{thm:main} with \(\varepsilon \approx 0.4\) and
\(\delta \approx 0.05\) predicts
\(W_2 \le C_1 \cdot 0.63 + C_2 \cdot 0.065 + C_3 \cdot 0.22\).
Corollary~\ref{cor:vr} translates this into a predicted variance gap;
the empirical \(|1 - \mathrm{vr}| = 0.07\) is consistent with the
bound under moderate values of \(C_1, C_3\).

\paragraph{What is and is not proved.}
Proposition~\ref{thm:main} bounds the law of the trained predictor
by training-error and discretisation contributions, under the
Gaussian-separable assumption (A1)--(A6), using the
kernel-smoothed periodogram loss with bandwidth
$h^\star = N_s^{-1/(2s_x+d)}$. The rigorous single-realisation
exponent is $\beta_{\rm sr} = s_x/(2(2s_x+d))$; the bias-only
exponent $\beta_{\rm bias} = \tfrac12\min(s_x/d,\,s_t)$
(Remark~\ref{rem:bias-only-rate})
is the faster rate observed empirically when the spectral target
is averaged across realisations, but is a heuristic statement
about the bias rather than a single-realisation guarantee. The
non-separable case is addressed by Proposition~\ref{thm:nonsep}
at the cost of the weaker rate $\beta' = s/(2(d+1))$.

The L\'evy-driven case is treated only partially:
Proposition~\ref{thm:levy-ext} bounds the Gaussian projection
(second-order structure) without Gaussianity but does not control
higher cumulants; Proposition~\ref{prop:levy-ito} extends to the
full law under (A7) at the cost of an additive $4\sqrt J$
correction, which is $O(\sqrt\nu)$ for VG drivers and is
non-vanishing as $N_s \to \infty$. The pathwise Lemma~\ref{lem:mm}
applied with a spline-interpolated L\'evy driver picks up an
unbounded-variation jump-component interpolation error that does
not vanish as $N_s \to \infty$ (see the remark following its proof);
this is the principal obstacle to extending the rigorous main theorem
to the L\'evy regime, and the open problem of a rate-form bound that
scales with the spectral training tolerance is acknowledged.
Proposition~\ref{cor:berry-esseen} quantifies cumulant-level
tightness of the Gaussian-projection approximation at the marginal.
Conjecture~\ref{prop:ntk} is conditional on the open positivity
$\lambda_{\min}(K^{\rm spec}) > 0$ (Remark~\ref{rem:lambda-min});
the end-to-end training-time statement is therefore not a theorem
but a conjecture.

The remaining open gaps are:
(i) global optimality of the empirical risk minimiser in the
non-convex landscape ($\varepsilon, \delta$ are assumed as attained
endpoint values, not shown to be attainable by any finite
computational procedure);
(ii) a fully rigorous full-law convergence theorem for L\'evy-driven
SPDEs that scales with the training loss rather than the
jump-energy $J$;
(iii) lower-bounding $\lambda_{\min}(K^{\rm spec}) > 0$ for the
log-PSD NTK, which would upgrade Conjecture~\ref{prop:ntk} to a
theorem;
(iv) bounding the spline-interpolation error of an
unbounded-variation L\'evy driver and absorbing it into a
constructive $\delta$.
Items (i)--(iv) are concrete future work.

\subsubsection{End-to-end training-time bound}
\label{subsec:training-time}

\begin{conjecture}[End-to-end training-time bound, conditional on
NTK positivity]
\label{prop:ntk}
\emph{Conditional on the open conjecture that the spectral-loss NTK
satisfies $\lambda_{\min}(K^{\rm spec}) > 0$
(Remark~\ref{rem:lambda-min}):} under (A1)--(A6) and (A9), with
initial spectral loss
$\varepsilon_0=\mathcal{L}_{\mathrm{spec}}(u_\theta(0))$:
\begin{equation}
\label{eq:w2-time}
W_2(T)\;\le\;
C_1\!\sqrt{\varepsilon_0\,e^{-\lambda_{\min}(K)\,T}
+\frac{C_{\mathrm{ntk}}}{\sqrt{m}}}
\;+\;C_2\,N_s^{-\beta_{\rm sr}}\;+\;C_3\sqrt\delta.
\end{equation}
To reach $W_2(T_\star)\le\eta$, provided
$\eta>C_2 N_s^{-\beta_{\rm sr}}+C_3\sqrt\delta
+C_1\sqrt{C_{\mathrm{ntk}}/\sqrt{m}}$,
it suffices to train for
\begin{equation}
\label{eq:training-time}
T_\star\;\ge\;
\frac{1}{\lambda_{\min}(K)}\,
\log\!\left(
\frac{C_1^2\,\varepsilon_0}
{\displaystyle\Bigl(\eta-C_2 N_s^{-\beta_{\rm sr}}-C_3\sqrt\delta\Bigr)^2
-\dfrac{C_1^2\,C_{\mathrm{ntk}}}{\sqrt{m}}}
\right).
\end{equation}
\end{conjecture}

\begin{proof}[Argument conditional on $\lambda_{\min}(K^{\rm spec}) > 0$]
Under (A9) and gradient flow, we assume the PL-type gradient
inequality
$\tfrac{d}{dT}\mathcal{L}_{\mathrm{spec}}
\le -\lambda_{\min}(K)\mathcal{L}_{\mathrm{spec}}$
holds for the log-PSD loss. This is the analogue of the standard
NTK decay established by~\cite{Wang2022} for $L^2$ losses; its
transfer to the non-linear, non-convex log-PSD loss is formal and
requires $\lambda_{\min}(K)>0$, which is open (Remark~\ref{rem:lambda-min}).
Granting the inequality, Gr\"onwall gives
$\mathcal{L}_{\mathrm{spec}}(u_\theta(T))
\le\varepsilon_0 e^{-\lambda_{\min}(K)T}+C_{\mathrm{ntk}}/\sqrt{m}$.
Substituting into Proposition~\ref{thm:main} yields~\eqref{eq:w2-time}.
Setting $W_2(T_\star)\le\eta$, isolating the exponential term, and
taking logarithms gives~\eqref{eq:training-time}.
$\hfill\square$
\end{proof}

\begin{remark}[Lower-bounding $\lambda_{\min}(K)$ for the spectral loss]
\label{rem:lambda-min}
Equation~\eqref{eq:training-time} is explicit in $\eta,N_s,\delta,m,\varepsilon_0$
but depends on $\lambda_{\min}(K)$.
For standard $L^2$ losses, $\lambda_{\min}(K)>0$ is guaranteed by the
positive-definiteness of the NTK in the infinite-width
limit~\cite{Wang2022}. For the log-PSD loss, the NTK Gram matrix
entries involve $\nabla_\theta[\log\widehat{S}^{xy}(u_\theta)](x_i)
\cdot\nabla_\theta[\log\widehat{S}^{xy}(u_\theta)](x_j)$,
differing from the regression kernel by a $1/\widehat{S}^{xy}$
preconditioning. Lower-bounding $\lambda_{\min}(K^{\mathrm{spec}})$
for this loss is the principal remaining gap in
making~\eqref{eq:training-time} fully constructive.
\end{remark}

\paragraph{Relation to existing PINN error analysis.}
The result is a strict extension of the Mishra--Molinaro
framework~\cite{MishraMolinaro2022}: removing the spectral loss from
training drops the bound on \(\varepsilon\), so only
\eqref{eq:pde-path} survives and the result reduces to the original
\(W_2 \le C_3\sqrt\delta\); removing the PDE residual drops the bound
on \(\delta\), so only \eqref{eq:spectral-path} survives, yielding a
pure spectral-matching bound which is novel. The two paths give
independent upper bounds on the same \(W_2\); they act on orthogonal
aspects of the law (\(\delta\) on pathwise \(L^2\) distance,
\(\varepsilon\) on the spectral measure) and the displayed
\(C_1\sqrt\varepsilon + C_2 N_s^{-\beta_{\rm sr}} + C_3\sqrt\delta\) is the
loose union via \(\min(a,b)\le a+b\), used here for symmetry with the
training objective rather than for tightness.


\section*{S2 \quad Numerical sanity check on Lemma~\ref{lem:single-realisation}}
\label{app:smoothed-verification}

We provide a small numerical sanity check that kernel-smoothing the
single-realisation periodogram reduces $L^2$ estimation error on a
synthetic Sobolev-regular spectral density, supporting the rigorous
appeal to Lemma~\ref{lem:single-realisation} in
Proposition~\ref{thm:main}.

\paragraph{Setup.} 1D torus with $N \in \{64, 128, 256, 512, 1024, 2048\}$
modes; spectral density $\mu(k) = (1 + k^2)^{-s}$ with $s = 1$
(matching the Variance-Gamma regularity in the experiments).
At each $N$ we sample $2048$ independent periodogram realisations,
each $\hat S(k) = \tfrac12(\chi^2_2)\cdot\mu(k)$ with chi-squared draws
indexed by $k$, then sweep a uniform-kernel bandwidth
$h \in \{0, 1, 2, \dots\}$ in cells and report the per-mode MSE
$\frac{1}{N}\sum_k(\hat S^{(h)}(k) - \mu(k))^2$ at the optimal $h$.
Implementation in
\texttt{scripts/verification/verify\_smoothed\_periodogram.py}.

\paragraph{Observed behaviour.}
At every $N$ the MSE-optimal smoothed estimator outperforms the raw
periodogram by a factor of $\sim\!0.45$ (i.e., smoothed MSE is roughly
half of raw MSE), with the optimal bandwidth $h^\star_{\rm cell} = 1$
(a 3-point uniform smoother) across the tested range. The constant
factor confirms the qualitative claim of
Lemma~\ref{lem:single-realisation}: smoothing strictly improves the
single-realisation $L^2$ estimation error.

\paragraph{Caveat on the rate.}
The predicted asymptotic rate
$N^{-s/(2s+d)} = N^{-1/3}$ (here $d = 1$) is not cleanly recovered in
this finite-$N$ regime: the per-mode-averaged MSE in our setup is
itself a discretisation-dependent functional whose decay reflects both
the estimator efficiency and the per-mode normalisation. A clean test
of the rate would require either fixing the reference frequency
window in continuum units (so the comparison normalisation does not
shift with $N$) or measuring pointwise error at a fixed $k$ rather
than the per-mode average. The estimator-improvement factor at the
optimal bandwidth is what
Lemma~\ref{lem:single-realisation} guarantees qualitatively; the
rate result is from the established
Dahlhaus--Brillinger non-parametric spectral estimation literature
\cite{Dahlhaus1996,Brillinger1981} and is invoked as such in the proof.

\paragraph{Connection to the Burgers experiments.}
The Burgers experiments in Section~\ref{subsec:variance-gap} of the
main paper use the raw periodogram in (A4) at training time, not the
smoothed periodogram. Their interpretation in light of the rigorous
single-realisation bound is therefore: the bias-only heuristic rate
$\beta_{\rm bias} = 1/4$ describes the empirical grid-convergence
behaviour observed there, while the rigorous single-realisation rate
$\beta_{\rm sr} = 1/8$ guarantees consistency once the loss is
evaluated against the kernel-smoothed periodogram instead.
Re-running the Burgers training with kernel-smoothed targets is an
expected operational confirmation but is not strictly required for
the theory to apply --- the rigorous bound is a statement about the
training procedure that uses the smoothed target. We do not retrain
in this revision; the existing empirical results stand as evidence
for the bias-only heuristic rate, which the theory's
Remark~\ref{rem:bias-only-rate} explicitly identifies.

\section*{S3 \quad L\'evy intensity threshold sweep}
\label{app:levy-tau}

We investigate whether a L\'evy intensity loss
\(\mathcal{L}_{\lambda}\), a soft count of field exceedances
above threshold \(\tau\sigma_{\mathrm{ref}}\), adds value
over cumulant matching alone, in particular at extreme thresholds
where cumulant matching might be expected to miss the rarest
events.

\subsection{Setup}

We train 18 models on the Heston-VG volatility channel with
randomised sub-grid: two loss modes
(\texttt{levy}: intensity only; \texttt{levy\_cumulant}: intensity
+ fourth-cumulant) crossed with three thresholds
\(\tau \in \{2, 3, 4\}\sigma_{\mathrm{ref}}\) and three training
seeds, each evaluated on 5 validation seeds (90 evaluations
total). The intensity loss weight is \(w_\lambda = 10{,}000\).
All other hyperparameters are as in Table~\ref{tab:hyper}.

\subsection{Results}

\begin{table}[h]
\centering
\caption{L\'evy intensity threshold sweep, V channel. Mean
$\pm$ standard deviation across 15 evaluations (3 train
$\times$ 5 val seeds). Reference: vr = 1.0,
$\kappa_4 = 27.3$, jf\(_3\sigma\) = 0.0238.}
\label{tab:levy-tau}
\begin{tabular}{l c c c c c}
\toprule
mode & $\tau$ & var.\ ratio & kurtosis & jf\(_3\sigma\) pred & jf\(_3\sigma\) ref \\
\midrule
levy & 2 & $0.910 \pm 0.095$ & $10.1 \pm 3.1$ & $0.0264$ & $0.0238$ \\
levy & 3 & $0.921 \pm 0.091$ & $8.1 \pm 2.2$ & $0.0267$ & $0.0238$ \\
levy & 4 & $0.930 \pm 0.099$ & $8.5 \pm 1.0$ & $0.0319$ & $0.0238$ \\
\midrule
levy\_cumulant & 2 & $0.923 \pm 0.098$ & $24.7 \pm 1.0$ & $0.0269$ & $0.0238$ \\
levy\_cumulant & 3 & $0.923 \pm 0.096$ & $25.1 \pm 0.7$ & $0.0256$ & $0.0238$ \\
levy\_cumulant & 4 & $0.921 \pm 0.092$ & $24.9 \pm 0.5$ & $0.0256$ & $0.0238$ \\
\bottomrule
\end{tabular}
\end{table}

\subsection{Conclusions}

Three conclusions follow from Table~\ref{tab:levy-tau}:

\begin{enumerate}
\item \textbf{L\'evy-only misses kurtosis at every threshold.}
Even at $\tau = 4\sigma$ (the most extreme threshold, targeted
at the rarest events), the \texttt{levy} variant reproduces
variance (\(\mathrm{vr} \approx 0.93\)) but achieves kurtosis
of only 8.5 vs.\ reference 27.3. The intensity loss, counting
exceedances above a threshold, does not carry the same information
as the fourth-cumulant penalty.

\item \textbf{The threshold $\tau$ is irrelevant for
\texttt{levy\_cumulant}.} Kurtosis, variance ratio, and jump
frequency are statistically identical across $\tau \in
\{2, 3, 4\}\sigma$ (within one standard deviation). The
cumulant term dominates, and the intensity loss adds nothing.

\item \textbf{L\'evy intensity matching is redundant.} The
\texttt{levy\_cumulant} results are indistinguishable from the
\texttt{cumulant}-only results in Table~\ref{tab:heston-2x4}.
We conclude that the L\'evy intensity penalty can be safely
dropped from the loss without any degradation in the reported
metrics.
\end{enumerate}

The \texttt{levy}-only variant reproduces variance correctly
(same as all randomised-subgrid variants) but leaves kurtosis
near-Gaussian, confirming the orthogonality of variance-gap
closure and kurtosis matching identified in
Section~\ref{subsec:heston-diag}.

% ─────────────────────────────────────────────────────────────────────────────
\section*{S4 \quad Full-law convergence for L\'evy drivers via Zolotarev
metrics}
\label{app:zolotarev}
% ─────────────────────────────────────────────────────────────────────────────

This appendix resolves the two obstacles identified in
Section~3 of the main text for the L\'evy extension: the
non-vanishing $4\sqrt J$ Gaussian-projection gap of
Proposition~\ref{prop:levy-ito}, and the $\Theta(1)$ spline-smoothing
jump floor of the pathwise leg (Remark after Lemma~\ref{lem:mm}). It
also provides the first finite law-distance bound for $\alpha$-stable
drivers with $\alpha < 2$, for which $W_2 = \infty$ and
Assumption~(A7) fails. The price is a change of metric: results are
stated in the Zolotarev ideal metric $\zeta_s$,
$s \in (1, 2]$~\cite{Zolotarev1976}, and transferred back to
Wasserstein distances where the latter are finite
\cite{Rio2009,BolbotowskiBouchitte2025}.

\subsection{Setting}

Let $H := L^2(\Omega \times [0,T])$ and let the driver be bandlimited
to modes $|k| \le K$, $J_K := \#\{k\}$, consistently with the sub-grid
reading of (A3). Write $\lambda_k$ for the symbol of the parabolic
operator on mode $k$, $g_k(\tau) = e^{-\lambda_k\tau}\mathbf 1_{\tau>0}$,
and $h$ for the temporal resolution of the driver grid. Throughout,
$h \le \lambda_K^{-1}$ (the FD solver's own stability regime). Driver
classes: \emph{(VG)} Variance-Gamma, jump second moment
$m_2 := \int z^2\nu(dz) < \infty$, $s = 2$; \emph{($\alpha$S)}
symmetric $\alpha$-stable, $\alpha \in (1,2)$, $s \in (1,\alpha)$,
compensated (It\^o) convention.

The experimental pipeline couples all objects on one probability
space: the FD reference $u_{\mathrm{FD}}$ is driven by a realisation
$L$; training evaluates residuals against the spline smoothing
$\tilde L = \mathcal S_h L$ of the \emph{same} increments. We use two
exact mild solutions as intermediaries: $\bar u$ (driven by
$\tilde L$) and $u^\star$ (driven by $L$). The error chain is
\[
u_\theta
\;\overset{\text{Lem.~\ref{lem:ls-stability}}}{\longleftrightarrow}\;
\bar u
\;\overset{\text{Lem.~\ref{lem:smoothing-law}}}{\longleftrightarrow}\;
u^\star
\;\overset{\varepsilon_{\mathrm{num}}}{\longleftrightarrow}\;
u_{\mathrm{FD}},
\]
where $\varepsilon_{\mathrm{num}}$ is the strong numerical error of
the reference solver, cited from the literature
\cite{PeszatZabczyk2007} and not re-derived here.

\subsection{The Zolotarev metric on $H$}

For $s = 1 + \beta$, $\beta \in (0,1]$, let
\[
\Lambda_s := \bigl\{f : H \to \mathbb R \ \text{Fr\'echet-}C^1,\
\|f'(x) - f'(y)\|_{H^*} \le \|x-y\|^\beta \ \forall x, y\bigr\},
\qquad
\zeta_s(\mu,\nu) := \sup_{f \in \Lambda_s}
\Bigl|\int f\,d\mu - \int f\,d\nu\Bigr|.
\]
$\zeta_s$ is \emph{ideal of order $s$}: translation-regular and
$s$-homogeneous under dilations~\cite{Zolotarev1976}. Since
$\Lambda_s$ contains every continuous linear functional,
$\zeta_s(\mu,\nu) < \infty$ requires equal means; we therefore state
bounds for the composite
\[
\rho_s(X, Y) := \zeta_s\bigl(\mathcal L(X - \mathbb EX),\,
\mathcal L(Y - \mathbb EY)\bigr) + \|\mathbb EX - \mathbb EY\|_H,
\]
which metrises the same topology (weak convergence plus $s$-moment
convergence). Two transfer results recover Wasserstein control:
$\tfrac14 W_2^2 \le \zeta_2$ on $\mathcal P_2(\mathbb R^d)$ for every
$d$, with sharp dimension-free
constant~\cite[Thm.~3.1]{BolbotowskiBouchitte2025}; and
$W_s^s \le c_s\,\zeta_s$ on $\mathbb R$ for
$s \in (1,2]$~\cite{Rio2009}.

\begin{lemma}[Taylor coupling bound]\label{lem:zeta-coupling}
Let $X, Y \in L^s(\mathbb P; H)$ with $\mathbb EX = \mathbb EY$. Then
\[
\zeta_s(\mathcal L X, \mathcal L Y) \;\le\;
\bigl(\mathbb E\|Y\|^s\bigr)^{\frac{s-1}{s}}
\bigl(\mathbb E\|X-Y\|^s\bigr)^{\frac 1s}
+ \mathbb E\|X-Y\|^s.
\]
\end{lemma}

\begin{proof}
Fix $f \in \Lambda_s$ and set
$g(x) := f(x) - f(0) - \langle f'(0), x\rangle$; then $g \in
\Lambda_s$, $g'(0) = 0$, and the affine part cancels in
$\mathbb E f(X) - \mathbb E f(Y)$ by mean equality. From $g'(0) = 0$
and the H\"older property, $\|g'(x)\| \le \|x\|^\beta$. Along the
segment from $Y$ to $X$,
\[
|g(X) - g(Y)| \le \int_0^1
\bigl(\|Y\| + t\|X-Y\|\bigr)^\beta dt\,\|X-Y\|
\le \|Y\|^{s-1}\|X-Y\| + \|X-Y\|^s,
\]
using $(a+b)^\beta \le a^\beta + b^\beta$ for $\beta \le 1$. Take
expectations and apply H\"older with exponents
$\tfrac{s}{s-1}, s$ to the first term.
\end{proof}

\begin{lemma}[Residual stability in $L^s$]\label{lem:ls-stability}
Linear case: pathwise, $u_\theta - \bar u$ solves the deterministic
parabolic equation with source equal to the training residual
$r_\theta$ (both objects see the same smoothed driver), so parabolic
stability gives
$\mathbb E\|u_\theta - \bar u\|_H^s \le C_{\mathrm{stab}}^s\,\delta_s$
with $\delta_s := \mathbb E(\|r_\theta\|_{H'} +
\|u_\theta(\cdot,0) - u_0\|)^s$. Semilinear case (Burgers-VG): same
statement under the well-posedness hypothesis of
\cite[Thm.~3.4]{MishraMolinaro2022}, with $C_{\mathrm{stab}}$ from a
Gr\"onwall argument on a ball containing both trajectories.
\end{lemma}

Note that $\delta_s$ is precisely a trainable residual loss with
exponent $s$; see Remark~\ref{rem:loss-exponent}.

\begin{lemma}[Driver smoothing vanishes in law]\label{lem:smoothing-law}
Let $Z := u^\star - \bar u = \int_0^\cdot S_{\cdot - \sigma}\,
d(L - \tilde L)(\sigma)$, with per-mode mismatch kernel
$R_k(t,\sigma) = g_k(t-\sigma) - (g_k * K_h)(t-\sigma)$, $K_h$ the
spline kernel of mass $1$ and support $O(h)$. Then, for
$h \le \lambda_K^{-1}$:
\[
\bigl(\mathbb E\|Z\|_H^s\bigr)^{1/s} \;\le\;
\begin{cases}
C_{\mathrm{VG}}\; J_K^{1/2}\, h^{1/2}, & \text{(VG)},\ s = 2,\\[3pt]
c_{\alpha,s}\, C\,(T J_K)^{1/\alpha}\, h^{1/2}, &
\text{($\alpha$S)},\ s < \alpha,
\end{cases}
\]
with $C_{\mathrm{VG}}^2 = c_0\, m_2 T$ explicit.
\end{lemma}

\begin{proof}
\emph{Kernel sections.} For fixed $\sigma$, the section
$t \mapsto R_k(t, \sigma)$ consists of a jump-mismatch region
$t \in (\sigma, \sigma + ch)$ where $|R_k| = O(1)$, and a smooth tail
where $|R_k| \le c\,h\,\lambda_k e^{-\lambda_k(t-\sigma)}$. Hence
\[
\|R_k(\cdot, \sigma)\|_{L^2_t}^2 \le c_0 h + c_1 h^2 \lambda_k
\le C^2 h
\quad\text{for } h\lambda_k \le 1.
\]

\emph{(VG).} The It\^o isometry for compensated Poisson integrals in
Hilbert space~\cite{PeszatZabczyk2007} gives
\[
\mathbb E\|Z\|_H^2 = \sum_{|k|\le K}\int_0^T
\|R_k(\cdot,\sigma)\|_{L^2_t}^2\, m_2\, d\sigma
\le m_2\, T\, J_K\, C^2 h,
\]
and $(\mathbb E\|Z\|^s)^{1/s} \le (\mathbb E\|Z\|^2)^{1/2}$ for
$s \le 2$.

\emph{($\alpha$S).} $Z$ is a jointly S$\alpha$S element of $H$ with
scale parameter
\[
\sigma_Z^\alpha = \sum_{|k| \le K}\int_0^T
\bigl\|R_k(\cdot, \sigma)\bigr\|_{L^2_t}^{\alpha}\, d\sigma
\]
(the $H$-norm is taken \emph{inside}, the $L^\alpha$ integral over
the random-measure variable $\sigma$ \emph{outside}). By the LePage
series representation and the equivalence of moments of stable norms
\cite[Ch.~3]{SamorodnitskyTaqqu1994}\cite[Ch.~5]{LedouxTalagrand1991},
$(\mathbb E\|Z\|_H^s)^{1/s} \le c_{\alpha,s}\,\sigma_Z$ for every
$s < \alpha$. Substituting the section estimate,
$\sigma_Z \le (T J_K)^{1/\alpha} C h^{1/2}$.
\end{proof}

\begin{remark}[Numerical verification]\label{rem:lemA3-verify}
A single-mode Monte-Carlo verification
(\texttt{scripts/verify\_lemma\_A3.py}; $T = 1$, fine grid $2^{16}$,
$\lambda \in \{4, 25\}$, $h \in [2.4\times 10^{-4},
6.3\times 10^{-2}]$, fits restricted to $\lambda h \le 1$) confirms
both rates. VG ($256$ replications): log--log slopes $+0.492$ and
$+0.504$ against the predicted $+1/2$, with the ratio
$\mathbb E\|Z\|^2/(m_2 T h)$ flat at $\approx 1/6$ across two decades
of $h$, suggesting the sharp constant $c_0 = 1/6$ for linear splines.
$\alpha$S ($\alpha = 1.5$, $1024$ replications): the robust median
estimator gives slopes $+0.495$ and $+0.483$, confirming $h^{1/2}$.
(An $L^{1.2}$-moment estimator is unreliable here — it has infinite
variance for $2s > \alpha$ — and is reported only as a consistency
check.) An earlier draft of this lemma claimed the rate
$h^{1/\alpha}$ for the stable case via an incorrect norm composition
in $\sigma_Z$; the verification falsified that exponent and the
statement above is the corrected one.
\end{remark}

\begin{lemma}[Cylindrical lift of the sharp Rio inequality]
\label{lem:bb-lift}
For centred $\mu, \nu \in \mathcal P_2(H)$, $H$ separable Hilbert:
$W_2^2(\mu,\nu) \le 4\,\zeta_2(\mu,\nu)$.
\end{lemma}

\begin{proof}
Let $\pi_m : H \to \mathbb R^m$ project onto the first $m$
coordinates of a fixed orthonormal basis. For
$f \in \Lambda_2(\mathbb R^m)$, the chain rule gives
$(f \circ \pi_m)'(x) = \pi_m^* f'(\pi_m x)$ and
$\|\pi_m^*(f'(\pi_m x) - f'(\pi_m y))\| \le \|\pi_m(x-y)\| \le
\|x - y\|$, so $f \circ \pi_m \in \Lambda_2(H)$ and
$\zeta_2(\pi_m\#\mu, \pi_m\#\nu) \le \zeta_2(\mu,\nu)$. By
\cite[Thm.~3.1]{BolbotowskiBouchitte2025} in $\mathbb R^m$ (constant
$4$, dimension-free),
$W_2^2(\pi_m\#\mu, \pi_m\#\nu) \le 4\zeta_2(\mu,\nu)$. Finally
$W_2(\pi_m\#\mu, \pi_m\#\nu) \uparrow W_2(\mu,\nu)$: optimal
couplings of the projected pairs are tight, any weak limit point is a
coupling of $(\mu,\nu)$, and the cost converges by monotone
convergence along the increasing projections.
\end{proof}

\subsection{Main result}

\begin{proposition}[Full-law convergence for L\'evy-driven Bochner
PINNs]\label{prop:zeta-main}
Assume (A1)--(A4), bandlimited drivers, the shared-driver coupling,
and Lemma~\ref{lem:ls-stability}'s stability hypothesis. Let $s = 2$
(VG) or $s \in (1, \alpha)$ ($\alpha$S), and set
$\Delta_s := (\mathbb E\|u_\theta - u_{\mathrm{FD}}\|_H^s)^{1/s}$.
Then
\[
\Delta_s \;\le\; C_{\mathrm{stab}}\,\delta_s^{1/s}
+ C_{\mathrm{sm}}\, J_K^{\gamma}\, h^{1/2}
+ \varepsilon_{\mathrm{num}},
\qquad \gamma = \tfrac12\ \text{(VG)},\ \tfrac1\alpha\ \text{($\alpha$S)},
\]
and consequently, with
$M_s := (\mathbb E\|u_{\mathrm{FD}}\|_H^s)^{1/s}$,
\[
\rho_s\bigl(u_\theta, u_{\mathrm{FD}}\bigr) \;\le\;
M_s^{\,s-1}\,\Delta_s + \Delta_s^{\,s} + \Delta_s.
\]
All constants are explicit and finite for both driver classes; no
Gaussianity is assumed, no additive jump-energy floor appears, and
the bound is finite for $\alpha$-stable drivers with
$s < \alpha$.
\end{proposition}

\begin{proof}
The triangle inequality in $L^s(\mathbb P; H)$ along the chain
$u_\theta \to \bar u \to u^\star \to u_{\mathrm{FD}}$, using
Lemma~\ref{lem:ls-stability}, Lemma~\ref{lem:smoothing-law}, and the
solver error, gives the bound on $\Delta_s$. Apply
Lemma~\ref{lem:zeta-coupling} to the centred pair (centring changes
$\Delta_s$ by at most a factor $2$, absorbed into constants); the
mean discrepancy is bounded by
$\|\mathbb E u_\theta - \mathbb E u_{\mathrm{FD}}\| \le \Delta_s$ by
Jensen.
\end{proof}

\begin{corollary}[VG: Wasserstein bound without the $4\sqrt J$
floor]\label{cor:vg-floorfree}
By Lemma~\ref{lem:bb-lift}, for the VG class,
\[
W_2\bigl(\mathcal L(u_\theta), \mathcal L(u_{\mathrm{FD}})\bigr)
\;\le\; 2\sqrt{\rho_2}
\;\lesssim\; 2 M_2^{1/2} \Delta_2^{1/2}
\quad (\Delta_2 \text{ small}).
\]
Unlike Proposition~\ref{prop:levy-ito}, this vanishes as training and
grids refine ($\delta_2 \to 0$, $h \to 0$): the $4\sqrt J$ term of
the L\'evy--It\^o coupling is replaced by the vanishing
$C\sqrt{J_K h}$. The two bounds are complementary; their minimum is
operative.
\end{corollary}

\begin{corollary}[$\alpha$-stable: first finite law bound]
\label{cor:stable}
For $\alpha \in (1, 2)$ and $s \in (1, \alpha)$,
$\rho_s(u_\theta, u_{\mathrm{FD}})$ is finite and converges to zero
under training and grid refinement. On mollified marginals
$\langle u, \varphi\rangle$, $\varphi \in H$, Rio's
inequality~\cite{Rio2009} transfers this to
$W_s \le (c_s \zeta_s)^{1/s}$.
\end{corollary}

\begin{remark}[Loss exponent for heavy tails: a caution]
\label{rem:loss-exponent}
The bound above is stated for $s < \alpha$, which is required only so
that the $s$-th moments in Lemma~\ref{lem:ls-stability} are finite. It
is tempting to read this as a training prescription---that heavy-tailed
drivers \emph{should} be fitted with an $L^s$-residual, $s < \alpha$,
since the MSE residual ($s = 2$) has an infinite-variance stochastic
gradient when the driver is redrawn at each step. We caution against
that reading (see Appendix~S5). The infinite-variance gradient arises
only under fresh-draw (amortised or operator) training; for the
single-realisation training used here the driver is a \emph{fixed
finite field}, the residual-gradient variance is finite for every
exponent, and $L^2$ is well posed. Where the concern does apply, an
$L^s$-residual with $s < \alpha$ down-weights precisely the
extreme-forcing residuals that carry the heavy tail, and so
\emph{suppresses} the tail it is meant to model. The exponent
$s < \alpha$ is therefore a device for stating a finite bound, not a
recommended training loss.
\end{remark}

\begin{remark}[Scope]\label{rem:zeta-scope}
The results of this appendix use the shared-driver coupling of the
experimental pipeline and therefore certify the trained surrogate
against the reference realisation's law. The ensemble setting (no
coupling) requires distributional losses beyond second order —
second-order information cannot pin a non-Gaussian law (polyspectral
phase freedom). A characteristic-functional probe loss with an
Esseen-type smoothing transfer to $\zeta_s$ is the natural candidate
and is left to future work. The spectral leg of
Proposition~\ref{thm:main} remains complementary: it is what excludes
the variance-collapse failure mode, which
Proposition~\ref{prop:zeta-main} alone does not address.
\end{remark}

\begin{remark}[Scope of the certificate for tail risk]
\label{rem:zeta-not-var}
Proposition~\ref{prop:zeta-main} certifies the law of $u_\theta$ in the
Zolotarev metric $\zeta_s$, which metrises weak convergence together
with $s$-th moment convergence. It does \emph{not}, on its own, control
tail quantiles (VaR/ES): $\zeta_s$ acts on the smooth test class
$\Lambda_s$, whereas a quantile at a fixed level is a non-smooth
functional. Indeed, by the Taylor-coupling bound
(Lemma~\ref{lem:zeta-coupling}) one may make $\zeta_s$ arbitrarily small
while two $p$-quantiles differ by a fixed gap, by transporting a small
mass across a low-density region at the quantile. A quantile transfer is
restored only under a density lower bound $f \ge c > 0$ on a
neighbourhood of the quantile, giving
$\lvert \Delta\,\mathrm{VaR}_p \rvert \lesssim \zeta_s^{1/s}/c$. The
certificate is thus a law-distance guarantee, and we make no claim of
tail-quantile accuracy for the surrogate absent such a regularity
assumption.
\end{remark}

% ─────────────────────────────────────────────────────────────────────────────
\section*{S5 \quad Numerical verification of the reference solver and standing assumptions}
\label{app:audit}
% ─────────────────────────────────────────────────────────────────────────────

This appendix verifies, against direct computation, the quantities the
main convergence bound cites or assumes: the finite-difference reference
error $\varepsilon_{\mathrm{num}}$, the separability assumption of
Proposition~\ref{thm:main}, the tightness and residual prescription of
the Zolotarev certificate, and floating-point adequacy. Scripts are
listed in \S S5.5.

\subsection{The reference-solver error $\varepsilon_{\mathrm{num}}$}
For the linear per-mode equation $du = -\lambda u\,dt + dL$ with
symmetric $\alpha$-stable $L$, the mild solution
$u(T) = \int_0^T e^{-\lambda(T-s)}\,dL(s)$ is symmetric $\alpha$-stable
with closed-form scale
$c_\star = \big((1-e^{-\alpha\lambda T})/(\alpha\lambda)\big)^{1/\alpha}$;
the Euler--Maruyama reference produces an iterate that is itself exactly
$S\alpha S$ with the \emph{same} $\alpha$, so $\varepsilon_{\mathrm{num}}$
reduces to a pure scale error. Measured, the temporal order is
$p \approx 1.00$ uniformly in $\alpha \in \{1.3,\dots,1.9\}$ and in the
mode rate $\lambda$, and the tail index is preserved to within $0.006$
($\hat\alpha = 1.896,\,1.693,\,1.506,\,1.300$ against
$1.9,\,1.7,\,1.5,\,1.3$). The second-order central spatial stencil gives
$O(\Delta x^2)$ error provided the grid over-resolves the driver band
(Nyquist $\ge 2K$); under broadband forcing it degrades, the
near-Nyquist modes being damped at $\approx 0.4\times$ the true rate.
The configurations of Section~4 satisfy $\Delta x \ll \sqrt{D\,T}$ with
effective band well below Nyquist, so the spatial error is sub-percent.
In the nonlinear (stochastic Burgers) case, Variance-Gamma forcing
converges at order $1$; $\alpha$-stable forcing exhibits a rare-event
blow-up (robust across advection schemes, non-monotone in the driver
scale), the numerical signature of the a priori $L^\infty$ bound that
unbounded jumps violate---the regime to which the semilinear certificate
of \S S4 restricts itself. Hence
$\varepsilon_{\mathrm{num}} = O(\Delta t) + O(\Delta x^2)$ throughout the
regime used here.

\subsection{The separability assumption, measured}
We quantify non-separability as the relative error of approximating the
joint two-dimensional power spectral density by the product of its
marginals. The analytic $z$-field (linear Ornstein--Uhlenbeck per mode,
Lorentzian joint spectrum) has index $0.30$ ($0.13$ best rank-one); the
empirical fields, as the loss sees them, are $0.116$, $0.049$, $0.0015$,
ordered by damping. A cost ablation (the paper loss versus the paper
loss augmented with a joint-PSD term, three seeds) shows the separable
loss already reproduces the true coupling (learned non-separability
index $0.111$ against true $0.107$) with a \emph{better} variance ratio
($0.944$ versus $0.920$); the joint term improves only the low-power
spectral bins. The PDE residual, not the spectral loss, supplies the
space--time coupling, so separability is a quantitatively good
approximation here, controlled by field damping.

\subsection{Zolotarev certificate: tightness and residual exponent}
At the operating point the bound is $\approx 12\%$ of the field scale,
its distinctive smoothing term $C_{\mathrm{sm}}(J_K h)^{1/2}$ only
$3.8\%$ (subdominant, on the verified constant of
Lemma~\ref{lem:zeta-coupling} and its companions): non-vacuous, with
total tightness gated by the training residual. Under $\alpha$-stable
forcing the $L^2$ (MSE) residual has infinite-variance stochastic
gradients---the sample gradient variance does not converge, and a single
jump remains $\sim 1\%$ of the whole gradient at $N = 10^6$---whereas the
$L^s$ residual ($s < \alpha$) has finite-variance gradients \emph{when the
driver is redrawn at each step}. This does not, however, warrant $L^s$
residual training in general, and we correct the prescription here. For the
single-realisation training used in this paper the driver is a \emph{fixed
finite field}, so the minibatch residual-gradient variance is finite for
every exponent and $L^2$ is well posed; the infinite-variance concern applies
only to amortised or operator learning that draws fresh realisations.
Moreover, $L^s$ with $s < \alpha$ down-weights the large residuals located at
the extreme-forcing points that carry the heavy tail, and therefore
\emph{suppresses} the tail it is meant to model (empirically, a variance ratio
of $0.26$ and a $99$th-percentile tail ratio of $0.37$, against $1.04$ and
$0.96$ for $L^2$). The $L^s$ residual is thus either unnecessary (single
realisation) or a gradient-versus-tail tradeoff (fresh draws), not a remedy.
The certificate's scope for tail risk is stated in
Remark~\ref{rem:zeta-not-var}.

\subsection{Floating-point adequacy}
The higher-order autograd residual was checked for \texttt{fp32}
contamination (catastrophic cancellation between $u_t$ and $D\,u_{xx}$).
On an exact solution, where the residual is identically zero, the
\texttt{fp32} output is bitwise $0$; on a trained network the
\texttt{fp32} numerical error is $1.2\times10^{-7}$, six orders of
magnitude below the residual the network actually reaches
($3.4\times10^{-2}$). \texttt{fp32} is therefore adequate---residual
faithfulness is limited by approximation capacity, not precision, and
the method requires no \texttt{fp64} hardware.

\subsection{Reproducibility}
NumPy/PyTorch scripts (\texttt{float64} linear algebra):
\texttt{fd\_convergence\_stable\{,\_spatial,\_nonlinear\}.py}
($\varepsilon_{\mathrm{num}}$); \texttt{separability\_index.py} and
\texttt{ablation\_separability.py} (separability);
\texttt{zolotarev\_\{bound\_tightness,grad\_variance,var\_transfer\}.py}
(certificate); and \texttt{fp32\_residual\_floor.py} (floating point).
All are included in the repository named in the ``Code and data
availability'' statement of the main text.

\bibliographystyle{elsarticle-num}
\bibliography{references}

\end{document}
